\documentclass[11pt]{article}
\usepackage[utf8]{inputenc}
\usepackage{amsmath,amssymb,amsfonts}
\usepackage{float}
\parindent 0px

\title{Mathematical Foundations of Image Synthesis}
\author{ Kazi Asef Kabir (2305084)}
\date{January 4,2026}

\begin{document}
\maketitle

\section{Introduction}
Hello World!!\\
This is \LaTeX\ document.\\
A rectangle has side lengths of $(x+1)$ and $(x+3)$. \\
The equation of area of the rectangle is $${A(x)=x^2+4x+3}$$.\\

\section{Mathematical notation}
Superscripts : 
$$2x^3$$
$$2x^{34}$$
$$2x^{3x+4}$$
$$2x^{3x^4+5}$$

Subscripts:
$$x_1$$
$$x_{12}$$
$$x_{{12}_2}$$
$$x_{1_{2_3}}$$
$$a_0,a_1,a_2,a_3,\ldots,a_{100}$$


Greek Letters:
$$\pi$$
$$\Pi$$
$$\alpha$$
$$A=\pi r^2$$

Trigonomertric Functions:
$$y=\sin x$$
$$y=\sin 2x$$
$$y={\sin}^2 2x$$
$$y=\csc \theta$$
$$y=\csc \Theta$$
$$y=\sin^{-1} x$$
$$y=\arcsin x$$

Log Functions:
$$y=\log x$$
$$y=\log_4 x$$
$$y=\ln ({x^3+5})$$

Roots:
$$\sqrt{x^3+5}$$
$$\sqrt[4]{3x^7+2x+7}$$
$$\sqrt{\sqrt[3]{1+\sqrt{x^2+2x+5}}}$$

Fractions:
$$\frac{6}{7}$$\\
About $\displaystyle \frac{2}{3}$ of the glass is full.\\[16pt]
About $\frac{2}{3}$ of the glass is full.\\

$$\dfrac{2}{3}$$\\[6pt]

$$\frac{\sqrt{x+1}}{\sqrt{x+2}}$$\\

$$\frac{1}{1+\frac{1}{\sqrt{x+2}}}$$

\section{Bracket,Tables and Arrays}
The distributive property states that $a(b+c)=ab+ac$ for all $a,b,c \in \mathbb{R}$.\\[6pt]
The equivalence class of $a$ is $[a]$.\\[6pt]
The set $A$ is defined to be $\{1,2,3\}$.\\[6pt]
The movie ticket costs $\$11.50$.\\[6pt]
$$2\left (\frac{1}{x^2-1}\right )$$\\

$$2\left \langle \frac{1}{x^2-1}\right \rangle $$
$$2\left | \frac{1}{x^2-1}\right | $$
$$\left. \frac{dy}{dx}\right|_{x=1}$$
$$\left(\frac{1}{1+\left(\frac{1}{1+x}\right)}\right)$$


Tables:\\[6pt]
\begin{tabular}{|c|c|c|c|c|c|}
    \hline
    $x$ & 1 & 2 & 3 & 4 & 5\\ \hline
    $f(x)$ & 10 & 11 & 12 & 13 &14\\ \hline   
\end{tabular}
\\[6pt]
\begin{table}[H]
    \centering
 \def\arraystretch{1.5}   
    \begin{tabular}{|c|c|c|c|c|c|}
    \hline
    $x$ & 1 & 2 & 3 & 4 & 5\\ \hline
    $f(x)$ & $\frac{1}{2}$ & 11 & 12 & 13 &14\\ \hline
\end{tabular}
\caption{Practice Tables}
\end{table}
\end{document}